\subsection{(U) 型膜与整体单值变换}\label{sec:18g}
与前面一样，令 $\symbfit{\gamma}$ 为消失路径的特选基。给每个 $\gamma_k$ 配上 $\widetilde E$ 中的 (B) 型等变膜 $\widetilde\Delta_k^\#$，其中分次任意选取。此外，设 $L^\#$ 为包含在 $E^{\mathrm{shrink}}$ 中的 $E$ 中正合 Lagrangian 膜，$\widetilde L^\#=\widetilde L_+^\#\cup\widetilde L_-^\#$ 为相应的 (U) 型等变膜。

\StatementNumber{15}\begin{lemma}\label{st:18.15}
在导出范畴 $D(\widetilde{\A})=H^0(\Tw\widetilde{\A})$ 中，有
\begin{equation}\label{eq:18.25}
T_{\widetilde\Delta_1^\#}\cdots T_{\widetilde\Delta_m^\#}(\widetilde L_+^\#)\cong S(\widetilde L_-^\#).
\end{equation}
\end{lemma}
\begin{proof}
若先忽略膜结构和平移 $S$，可作如下简化解释。由\stref{推论}{17.17}，\eref{18.25}左端的对象同构于 Lagrangian 子流形 $\tau_{\widetilde\Delta_1}\cdots\tau_{\widetilde\Delta_m}(\widetilde L_+)$ 所给出的对象。\stref{引理}{18.2}说明，在带标架同痕意义下，$\widetilde\Delta_k$ 是将 $\gamma_k$ 看成 $\Psi:\Efamily\to\mathcal W$ 的消失路径时所对应的消失圈，因此它们的 Dehn 扭转是 $\Psi$ 的单值变换。\znote{284}{这里藏着一个需要点破的等式：为什么若干 Dehn 扭转的复合会同痕于沿圆周的单值变换？因为按 Picard--Lefschetz 公式（§16d），沿 $\lambda$ 的单值变换同痕于以各消失圈为核的 Dehn 扭转、按临界值沿路径出现的先后排成的乘积；而引理 18.2 保证 $\widetilde\Delta_k$ 确实是 $\Psi$ 的消失圈，故这些扭转就是 $\Psi$ 自身的扭转。次序（先 $\widetilde\Delta_1$ 后 $\widetilde\Delta_2$……）与沿 $\lambda$ 遇到临界值的先后一致，不能交换；换成逆序得到的是逆单值变换，这也解释了\eref{18.25}左端乘积为何按 $1,\ldots,m$ 排列。[来源：原书 §16d、§18g（印刷页 284）]}于是，全部 $\tau_{\widetilde\Delta_k}$ 的复合在 $\operatorname{Aut}(\widetilde E)$ 中同痕于 $\Psi$ 沿 $\mathcal W$ 中半径 $1-3\delta$ 的圆的单值变换。由\stref{引理}{18.1}，该单值变换在子集 $\widetilde E^{\mathrm{shrink}}$ 上的限制等于 $\iota_{\widetilde E}$；按定义，该子集包含 $\widetilde L_\pm$。因此，它将 $\widetilde L_+$ 映到 $\widetilde L_-$。

为得到完整陈述，需作更仔细的论证。$T_{\widetilde\Delta_k^\#}$ 在对象上的作用对应于带典范分次的 $\tau_{\widetilde\Delta_k}$ 的作用；由 Picard--Lefschetz 公式的分次版本，这也是自然赋予 $\Psi$ 的相应单值变换的分次。由\stref{引理}{18.6}，$\widetilde L_+^\#=(\widetilde L_+,\widetilde\alpha_+^\#,P_+^\#)$ 在此单值变换下的像是
\[
\bigl(\iota_{\widetilde E}(\widetilde L_+),\,
\iota_{\widetilde E}^*\widetilde\alpha_+^\#-1,\,
\iota_{\widetilde E}^*P_+^\#\bigr).
\]
它与 $\widetilde L_-^\#$ 相差分次中的常数 $-1$，以及 Pin 结构上的 $\otimes\lambda^{\mathrm{top}}(T\widetilde L_-)$。这一差异恰好就是\secref{12a}中定义的 Lagrangian 膜平移运算 $S$ 的作用；在\secref{12h}中，我们已验证它确实拟表示抽象的代数平移。
\end{proof}

\StatementNumber{16}\begin{note}\label{st:18.16}
若将 $S$ 看成使 Lagrangian 子流形的定向反向的运算，\eref{18.25}中出现 $S$ 似乎令人困惑。事实上，
\SourcePage{292}
相反的定向已经包含在 $\widetilde L_\pm^\#$ 的定义中，见\stref{笔记}{18.8}；平移仅仅恢复一个共同的选择。
\end{note}

\StatementNumber{17}\begin{proposition}\label{st:18.17}
$\widetilde L^\#$ 位于 $\Tw\F(\pi)$ 中由 $\widetilde\Delta_1^\#,\ldots,\widetilde\Delta_m^\#$ 生成的三角子范畴中。
\end{proposition}
\begin{proof}
这是\stref{引理}{6.3}的直接应用。$\widetilde L^\#$ 是等变膜，因而是 $\widetilde{\A}$ 中在对合 $\eps_{\widetilde{\A}}$ 下不变的对象。按定义，它同构于 $\widetilde L_+^\#$ 与 $\widetilde L_-^\#=\eps_{\widetilde{\A}}(\widetilde L_+^\#)$ 的直和。此外，$HF^*(\widetilde L_+^\#,\widetilde L_-^\#)=0$，因为底层 Lagrangian 子流形不相交。各 $\widetilde\Delta_k^\#$ 同样是 $\eps_{\widetilde{\A}}$-不变对象，并且由\eref{18.23}，对所有 $k<l$，$HF^*(\widetilde\Delta_k^\#,\widetilde\Delta_l^\#)$ 都是 $\Z/2$-不变的。\stref{引理}{18.15}恰好给出所需的关系\eref{6.2}，而按定义，$\widetilde{\A}^{\mathrm{inv}}=\F(\pi)$。
\end{proof}
